\section{Εκπαίδευση}
\label{Training}

Στο κεφάλαιο αυτό παρουσιάζεται ο τρόπος με τον οποίο εκπαιδεύονται τα μοντέλα που περιγράφηκαν στο κεφάλαιο \ref{Methodology}.

\subsection{Προετοιμασία των Δεδομένων Εκπαίδευσης}
\label{TrainingData}

\subsubsection{Εξαγωγή φωνημάτων για κάθε frame}
\label{PhonemeExtraction}
Για να έχει νόημα η εφαρμογή του όρου cross entropy, είναι απαραίτητο να δημιουργηθούν labels για κάθε frame των βίντεο. Αυτό επιτυγχάνεται με συνδυασμό της γνώσης του περιορισμένου περιβάλλοντος και μεθόδων αναγνώρισης φωνημάτων από ηχητικά σήματα ομιλίας. Να σημειωθεί ότι η παρακάτω διαδικασία μπορεί κάλλιστα να εφαρμοστεί και σε μη περιορισμένο περιβάλλον, αλλά οι επιπλέον περιορισμοί ενός περιορισμένου περιβάλλοντος επιτρέπουν απλοποιήσεις και διορθώσεις σε περίπτωση λανθασμένης προφοράς από τον ομιλητή.

Θεωρώντας γνωστό το φραστικό περιεχόμενο κάθε βίντεο (π.χ. στο GRID dataset γίνεται γνωστή η φράση που εκφέρεται από τον ομιλητή μέσω του τίτλου του αρχείου), μετατρέπεται η φράση σε αλληλουχία φωνημάτων μέσω του eSpeak-NG, δίχως ακόμα να υπάρχει ωστόσο γνώση για τον χρονισμό του κάθε φωνήματος. Χρησιμοποιείται ελαφριά κανονικοποίηση στα φωνήματα που εξάγονται ώστε το ``αλφάβητο'' των τελικών αλληλουχιών να ταυτίζεται με αυτό του μοντέλου wav2vec2, που χρησιμοποιείται για την εξαγωγή φωνημάτων από το ηχητικό περιεχόμενο. Το CTC μοντέλο του wav2vec2 εξάγει πιθανότητες κάθε φωνήματος για κάθε frame και, συγκρίνοντας με την αλληλουχία φωνημάτων που γνωρίζουμε ότι αναμένεται, παράγει μία τελική πρόβλεψη για την αρχή και το τέλος του κάθε φωνήματος. Μέσω αυτών των πληροφοριών, είναι δυνατόν να ολοκληρωθεί το labelling κάθε frame με το αντίστοιχο φώνημα ή το blank token, εφόσον το frame αντιστοιχεί σε χρονική στιγμή εκτός των φωνημάτων. Έτσι, πέραν της εξαγωγής των απαιτούμενων labels, αποφεύγονται τυχόντα προβλήματα από αστοχία κατά την προφορά της κάθε φράσης.

Η επεξεργασία των φωνημάτων δεν ολοκληρώνεται εκεί. Συγκεκριμένα, δημιουργείται ένα viseme grouping (ομαδοποίηση, δηλαδή, των όμοιων οπτικά φωνημάτων) με βάση τα χαρακτηριστικά των φωνημάτων. Για αναλυτικότερες πληροφορίες ως προς την ομαδοποίηση των φωνημάτων, ανατρέξτε στην ενότητα \ref{Phonemes}. Από τα hard targets που υπάρχουν ήδη για κάθε frame των βίντεο, δημιουργούνται soft targets με βάση τις παραπάνω ομάδες. Συγκεκριμένα, το 85\% αποδίδεται στο πραγματικό φώνημα, ενώ το υπόλοιπο 15\% μοιράζεται στα όμοια φωνήματα του πραγματικού με βάρη που ανταποκρίνονται στην ομοιότητα των φωνημάτων. Για ίδιο viseme, χρησιμοποιείται βάρος 0.95, ενώ για διαφορετικό υπολογίζεται similarity score ως εξής: κάθε φώνημα έχει τον τόπο, τον τρόπο άρθρωσης και το αν είναι έμφωνο ή όχι. Με βάση το πόσα κοινά features άρθρωσης έχουν δύο φωνήματα που δεν ανήκουν στο ίδιο viseme, αποδίδεται μεγαλύτερο βάρος σε αυτά με περισσότερες ομοιότητες και μικρότερο σε αυτά με τις λιγότερες. Έτσι, υπολογίζεται η αρθρωτική ομοιότητα (articulatory similarity) δύο φωνημάτων $i$ και $j$ ως εξής: $s(i,j) = 1 - \delta / 3$, όπου $\delta$ είναι ο αριθμός διαφορετικών features. Ο τελικός πίνακας ομοιότητας (similarity matrix) $A$ χτίζεται με τον ακόλουθο κανόνα:
$A_{ij} = \begin{cases}
1 & i = j \\
0 & \text{είτε το }i \text{ είτε το }j \text{ είναι blank token και } i \neq j\\
0.95 & \text{ίδια ομάδα viseme} \\
0.4 \cdot s(i,j) & \text{σε άλλη περίπτωση} 
\end{cases}
$

Τελικά τα soft targets χτίζονται με το 85\% στο πραγματικό φώνημα και το υπόλοιπο 15\% μοιρασμένο στα όμοια φωνήματα σε αναλογία με το similarity score των φωνημάτων, εκτός από την περίπτωση του blank token, κατά την οποία διατηρούνται τα hard targets. Αυτό συμβαίνει ώστε να ενισχυθεί η ικανότητα του μοντέλου να προβλέπει ακριβώς πότε λέγεται ένα φώνημα και πότε όχι, με αποτέλεσμα οι χρονισμοί των προβλέψεων να είναι ορθοί.

\paragraph{Μορφή των δεδομένων εκπαίδευσης}

Όπως παρουσιάστηκε παραπάνω, κάθε αντικείμενο του training (κάθε βίντεο) συμπληρώνεται από μία προβλεπόμενη αλληλουχία φωνημάτων σε επίπεδο frame. Για κάθε frame υπάρχει η πρόβλεψη του φωνήματος (κενό αν δεν εκφέρεται φώνημα), ενώ συμπληρωματικά υπάρχει και η συμπτυγμένη (collapsed) ακολουθία, στην οποία συνεχόμενα επαναλαμβανόμενα φωνήματα και κενά παραλείπονται και η οποία απαιτείται για τον όρο CTC της αντικειμεντικής συνάρτησης. Τέλος, συμπεριλαμβάνεται στις γνωστές πληροφορίες και ένα log-probability confidence score για κάθε πρόβλεψη, που έχει παραχθεί κατά τη δημιουργία των προβλεπόμενων αλληλουχιών. Τα log-probabilities περιορίζονται στο σύνολο $[-10,0]$ και έπειτα μεταφέρονται γραμμικά στο σύνολο $[0.5,1.5]$, με frames που δεν καλύπτονται από το span κάποιου φωνήματος (κενά frames) να παίρνουν το default βάρος 1. Τα βάρη αυτά επηρεάζουν, έπειτα, τη συνεισφορά του κάθε λάθους σε επίπεδο frame στον όρο cross entropy, ώστε το μοντέλο να τιμωρείται λιγότερο στις περιπτώσεις που η ίδια η πρόβλεψη του target είναι αβέβαιη.

\subsubsection{Data Augmentation}

Για να εμπλουτιστεί ιδιαίτερα κάποιο περιορισμένο σε μέγεθος dataset, κρίνεται αναγκαίο να εφαρμοστούν μέθοδοι data augmentation, δηλαδή επαύξησης δεδομένων. Αυτές οι μέθοδοι, εκμεταλλευόμενες τη γνώση μας στο αντικείμενο, παράγουν νέα δεδομένα από τα υπάρχοντα και εμπλουτίζουν τη μάθηση του μοντέλου. Είναι σημαντικό να τονιστεί ότι οι τεχνικές data augmentation που εφαρμόζονται πρέπει να συμβαδίζουν με τους στόχους μάθησης του μοντέλου. Για παράδειγμα, μία πιθανή μέθοδος data augmentation στην περίπτωση που τα δεδομένα βρίσκονται σε μορφή βίντεο είναι η αναστροφή της αλληλουχίας των frames. Ωστόσο, όταν το μοντέλο πρόκειται να μάθει σε κάποιο παρακείμενο στο οποίο έχει σημασία η σωστή χρονική αλληλουχία (όπως για παράδειγμα στο lip reading, όπου οι κινήσεις των χειλιών για την παραγωγή ενός φωνήματος είναι ευαίσθητες στη σωστή χρονική σειρά), δε θα μπορούσε να εφαρμοστεί η παραπάνω μέθοδος.

Οι τεχνικές data augmentation που τελικά εφαρμόζονται, λόγω της ορθότητάς τους για τον τελικό στόχο και λόγω της ταχείας εφαρμογής τους σε πολλαπλά δεδομένα, είναι ο καθρεπτισμός γύρω από τον κατακόρυφο άξονα (horizontal flip), η χρωματική αλλαγή (color jitter) και η αφαίρεση/επικάλυψη (cutout/occlusion). Σε κάθε περίπτωση οι μετασχηματισμοί αυτοί εφαρμόζονται σε ολόκληρο το βίντεο και όχι σε μεμονωμένα frames, ώστε να αποφευχθούν ανεπιθύμητα φαινόμενα αστάθειας. Ο καθρεπτισμός εφαρμόζεται με πιθανότητα $p=0.5$, ενώ η χρωματική αλλαγή πραγματοποιείται μέσω των παραγόντων της φωτεινότητας (brightness) με συντελεστή από 0.8 έως 1.2, της αντίθεσης (contrast) με παράγοντα 0.8-1.2 και του κορεσμού (saturation) με παράγοντα 0.9-1.1 και με πιθανότητα $p=0.4$. Τέλος, με τη μέθοδο της επικάλυψης, με πιθανότητα $p=0.25$ επιλέγεται τυχαία μία ορθογώνια περιοχή (8-20\% του ύψους και του πλάτους του frame) και αποκρύπτεται για όλα τα frame του δεδομένου βίντεο, προσομοιώνοντας την ύπαρξη ενός αντικειμένου που εμποδίζει τη θέαση συγκεκριμένης περιοχής. Αυτό έχει ως αποτέλεσμα η εκπαίδευση να γίνεται πιο έυρωστη. Να σημειωθεί ότι η διαδικασία του data augmentation εφαρμόζεται αποκλειστικά στο training split και όχι σε αυτά του validation και του test. Παρακάτω παρουσιάζονται μερικά παραδείγματα των παραπάνω μεθόδων επαύξησης δεδομένων.

\begin{figure}[H]
	\centering
	\includegraphics[width=0.7\textwidth]{horizontal_flip.png}
	\caption{Παρουσίαση του καθρεπτισμού (horizontal flip) για μερικά frames}
\end{figure}

\begin{figure}[H]
	\centering
	\includegraphics[width=0.7\textwidth]{color_jitter.png}
	\caption{Παρουσίαση της αλλαγής χρώματος (color jitter) για μερικά frames}
\end{figure}

\begin{figure}[H]
	\centering
	\includegraphics[width=0.7\textwidth]{random_cutout.png}
	\caption{Παρουσίαση της αφαίρεσης/επικάλυψης (cutout/occlusion) για μερικά frames}
\end{figure}

\subsubsection{Επιπλέον βήματα για την εκπαίδευση του audiovisual μοντέλου}
\label{AudiovisualPreprocessing}

Σε κάθε τμήμα ήχου προστίθεται θόρυβος και το αντίστοιχο τμήμα βίντεο υφίσταται παραμόρφωση, παράγοντας μία θορυβώδη εκδοχή του αρχείου ήχου, μία παραμορφωμένη εκδοχή του αρχείου βίντεο και δύο πίνακες αξιοπιστίας $r_a(t), r_v(t)$, οι οποίοι είναι ανεξάρτητοι για κάθε τροπικότητα και καταγράφουν την αξιοπιστία της εκάστοτε τροπικότητας για κάθε frame (πόσο αναμένεται να επηρεάστηκε η ποιότητα της εκάστοτε τροπικότητας από τον θόρυβο), λαμβάνοντας μέγιστη τιμή 1 (πλήρως αξιόπιστο frame) και ελάχιστη 0 (καθόλου αξιόπιστο). Έπειτα, από το αρχείο ήχου παράγεται το Mel-Spectrogram που αντιστοιχεί σε αυτό, το οποίο είναι η απαιτούμενη μορφή δεδομένων εισόδου για την τροπικότητα του ήχου και το προτεινόμενο μοντέλο.  Κάθε training instance για το audiovisual μοντέλο συμπληρώνεται από την πραγματική (ground truth) αλληλουχία φωνημάτων για κάθε frame, με την ίδια μορφή που χρησιμοποιείται στην περίπτωση του βασικού προτεινόμενου μοντέλου.

Είναι σημαντικό να τονιστεί ότι, κατά τη διαδικασία της εκπαίδευσης του audiovisual μοντέλου, απαιτείται η ύπαρξη θορύβου σε κάθε clip, τόσο στον ήχο του όσο και στο αντίστοιχο βίντεο. Η επιλογή αυτή είναι συνειδητή, ώστε να αποφευχθούν περιπτώσεις όπου το audiovisual μοντέλο μαθαίνει ότι μία από τις δύο τροπικότητες παρουσιάζει στατιστικά μικρότερη επιρροή από τον θόρυβο, αντί να μαθαίνει αξιοποιώντας και τα reliability σήματα $r_a(t)$ και $r_v(t)$. Εξαναγκάζοντας την ύπαρξη θορύβου και στα δύο streams δεδομένων, το μοντέλο μαθαίνει αναγκαστικά να συνδυάζει τα στοιχεία από τις δύο τροπικότητες και τελικά προσομοιώνει καλύτερα το περιβάλλον ύπαρξης θορύβου (και στον ήχο, π.χ. από ταυτόχρονη ομιλία άλλων ανθρώπων, και στο βίντεο, π.χ. με μερική ή πλήρη επικάλυψη των χειλιών από εμπόδια ή τεχνικά προβλήματα κατά την καταγραφή του βίντεο). Στη συνέχεια της ενότητας αυτής, παρουσιάζονται αναλυτικότερα τα βήματα προεπεξεργασίας που απαιτούνται.

\paragraph{Παραμόρφωση του βίντεο}

Σε αντίθεση με τον ήχο, στο βίντεο δεν προστίθεται θόρυβος με την αυστηρή, προσθετική έννοια, αλλά τα επιλεγμένα frames υφίστανται παραμόρφωση (μηδενισμό, πάγωμα ή θόλωμα, όπως περιγράφεται παρακάτω). Για λόγους συντομίας, στο υπόλοιπο κείμενο ο όρος ``θόρυβος'' στο βίντεο αναφέρεται σε αυτές τις παραμορφώσεις. Για κάθε βίντεο, ακριβώς μία αλληλουχία frames επηρεάζεται από την παραμόρφωση, και όχι διάφορα διάσπαρτα τμήματα. Αυτό αντικατοπτρίζει πραγματικά σενάρια που μπορεί να προκύψουν κατά την καταγραφή (ή τη ζωντανή μετάδοση) βίντεο, όπως τοποθέτηση εμποδίων μπροστά από τα χείλη (π.χ. χεριού ή μικροφώνου) ή τεχνικό πρόβλημα της κάμερας (π.χ. να χαθεί το focus και να εμφανιστεί θόλωση στο βίντεο), τα οποία είναι συνήθως ενιαία και σχετικά μεγάλου μήκους (στη γειτονιά του ενός δευτερολέπτου), αντί να είναι διάσπαρτα και σχετικά μικρού μήκους (πιο κοντά σε μεμονωμένα frames).

Αρχικά, επιλέγεται το span των frames του βίντεο στα οποία θα εφαρμοστεί παραμόρφωση. Το μήκος του συγκεκριμένου span επιλέγεται ομοιόμορφα στο range $[L_v - J_v, L_v + J_v]$, όπου $L_v$ είναι το κεντρικό (επιθυμητό) μήκος και $J_v$ η μέγιστη απόκλιση (jitter) από αυτό, με τις default τιμές $L_v = 30$ και $J_v = 8$. Σε αρχεία βίντεο στα 25fps, όπως για παράδειγμα αυτά του GRID, αυτό αντιστοιχεί σε θόρυβο στην περιοχή του ενός δευτερολέπτου (συνήθως λίγο παραπάνω από αυτό). Το span τοποθετείται τυχαία και ομοιόμορφα σε κάποια αρχική θέση εντός του clip, προσδίδοντας τυχαιότητα τόσο στην ακριβή θέση του θορύβου όσο και στο πόσο αισθητή είναι η παρουσία του (κρατώντας αυτήν την τιμή κοντά σε μία κεντρική επιθυμητή τιμή) ανάμεσα στα διάφορα clips.

Αφού έχουν επιλεχθεί με την παραπάνω μεθοδολογία τα frames τα οποία θα επηρεαστούν, εφαρμόζονται οι εξής μέθοδοι παραμόρφωσης, με τυχαίο τρόπο:

\begin{itemize}
\item zero (μηδενισμός): οι τιμές των pixels για τα επιλεγμένα frames τίθενται στην τιμή 0, προσομοιώνοντας το σενάριο όπου χάνεται πλήρως η πληροφορία των χειλιών (π.χ. λόγω της εμφάνισης εμποδίου μπροστά από τα χείλη). Για αυτόν τον τύπο παραμόρφωσης, η τιμή της αξιοπιστίας είναι 0, καθώς δεν υπάρχει οπτική πληροφορία που μπορεί να αξιοποιηθεί.

\item freeze (``πάγωμα''): κάθε frame του επιλεγμένου span αντικαθίσταται από το τελευταίο frame πριν από το span (ή, σε περίπτωση που το span ξεκινάει από το 1ο frame του βίντεο, το 1ο του span), προσομοιώνοντας ένα stream βίντεο το οποίο έχει παγώσει. Και πάλι το σκορ αξιοπιστίας για τα frames αυτά τίθεται ίσο με 0, μια και η οπτική πληροφορία που υπάρχει δεν μπορεί να χρησιμοποιηθεί για την εξαγωγή συμπερασμάτων όσον αφορά την έκφανση φωνημάτων. 

\item blur (θόλωμα): κάθε frame του span αντικαθίσταται από το αποτέλεσμα της συνέλιξης αυτού με ένα 25x25 Gaussian kernel (με σ=12), προσομοιώνοντας την απώλεια focus ή την εμφάνιση motion blur. Στη συγκεκριμένη περίπτωση, καθώς υπάρχει ένα ίχνος οπτικής πληροφορίας για την κίνηση των χειλιών που επιβιώνει, η τιμή της αξιοπιστίας για τα επηρεασμένα frames τίθεται ίση με 0.3 αντί για 0.
\end{itemize}

Από τη διαδικασία αυτή παράγεται το παραμορφωμένο βίντεο το οποίο αποτελεί είσοδο για το μοντέλο, όπως και το σήμα αξιοπιστίας $r_v(t), t = 0,1,...,T$, με τιμές 1 για κάθε μη επηρεασμένο frame $t$ και τιμές που ορίζονται από τα παραπάνω για τα παραμορφωμένα frames.

\paragraph{Προσθήκη θορύβου στον ήχο}

 Η προσθήκη θορύβου στον ήχο συμβαίνει σε 2 διαφορετικές αλλεπάλληλες στρώσεις, μία στρώση προσθετικού θορύβου που εφαρμόζεται πάντα και μία προαιρετική στρώση ολικού dropout του ηχητικού σήματος.
 
 Αρχικά, εφαρμόζεται ο θόρυβος προσθετικής φύσεως στο ηχητικό σήμα. Υπάρχουν δύο τύποι ηχητικού θορύβου, οι οποίοι επιλέγονται τυχαία για κάθε clip:
 
 \begin{itemize}
 \item white (λευκός θόρυβος): κλασικός i.i.d (independent identically distributed) γκαουσιανός θόρυβος
 
 \item babble (``βαβούρα''): μία πολύ γενική προσέγγιση βαβούρας από πολλαπλούς άλλους ομιλητές. Δημιουργείται παράγοντας λευκό θόρυβο, υπολογίζοντας το άθροισμά του (cumulative sum, προσομοιώνοντας έναν ολοκληρωτή και ενισχύοντας τις χαμηλές συχνότητες) και παίρνοντας την πρώτη διαφορά του. Έτσι, προσδίδεται στον θόρυβο spectral coloring, με αποτέλεσμα να παράγεται θόρυβος με διαφορετικά φασματικά χαρακτηριστικά από τον λευκό, ο οποίος χονδρικά προσομοιώνει την παρουσία πολλών διαφορετικών ομιλητών, χωρίς να χρειάζεται πρόσβαση σε κάποιο αρχείο βαβούρας.
 \end{itemize}
 
 Σε γενικές γραμμές, ο προσθετικός θόρυβος του ήχου θεωρήθηκε συνεχής και όχι διακοπτόμενος. Επιλέγεται τυχαία και ομοιόμορφα μία τιμή signal-to-noise ratio (SNR) $\rho$ σε dB από ένα ορισμένο range (με default τιμή $[-5,15]$ dB). Ο θόρυβος κανονικοποιείται ώστε να πετυχαίνεται το επιθυμητό SNR:
 
 \[
 n_s = n \cdot \sqrt{\frac{P_s/10^{\rho/10}}{P_n}}
 \]
 
 όπου $n$ είναι ο αρχικός θόρυβος, $n_s$ ο κλιμακωμένος θόρυβος και $P_s$, $P_n$ η ισχύς του σήματος και του θορύβου αντίστοιχα. Το τελικό σήμα μετά την προσθήκη του θορύβου επιδέχεται clipping στο σύνολο $[-1,1]$. Ο μηχανισμός αυτός εγγυάται την ύπαρξη θορύβου (έστω μικρού) σε κάθε audio clip.
 
 Στο δεύτερο επίπεδο, προστίθεται με πιθανότητα $p_a$ (με προεπιλεγμένη τιμή 0.3) μία περαιτέρω μορφή θορύβου. Το ηχητικό που αντιστοιχεί σε έναν ακριβή αριθμό frames $L_a$ (με default τιμή 15) αντικαθίσταται από το μηδενικό σήμα. Το φαινόμενο αυτό μοιράζεται σε 1 μέχρι $K_a$ (default τιμή 1, αλλά δίνεται από τον κώδικα ως επιλογή) διαφορετικά spans, προσομοιώνοντας μικρές αποκοπές του ηχητικού σήματος (π.χ. κόψιμο του μικροφώνου). 
 
 Όσον αφορά το σήμα αξιοπιστίας του ήχου $r_a(t)$, υπολογίζεται εφαρμόζοντας clipping του επιλεγμένου SNR:
 
 \[
 r_a = \operatorname{clip}\left(\frac{\rho - (-5)}{20 - (-5)}, 0, 1\right)
 \]
 
 όπου οι τιμές -5dB και 20dB επιλέχθηκαν ως τα όρια ενός αξιόπιστου ηχητικού σήματος: οτιδήποτε με SNR κάτω από -5dB θεωρείται πλήρως αναξιόπιστο, αντίθετα οτιδήποτε με SNR πάνω από 20dB θεωρείται πλήρως αξιόπιστο. Για frames στα οποία έχει εφαρμοστεί ολικό dropout του ηχητικού σήματος, η αξιοπιστία παίρνει την τιμή 0.
 
 \paragraph{Υπολογισμός των Mel-Spectrograms}
 
 Πριν την εκπαίδευση, το ηχητικό σήμα (θορυβώδες ή όχι), το οποίο αντιστοιχεί στους στόχους της εκπαίδευσης και σε κάθε τμήμα βίντεο, μετατρέπεται σε Mel-Spectrogram (βλ. ενότητα \ref{MelSpectrogram}). Καθώς ο υπολογισμός αυτός είναι πλήρως ντετερμινιστική διαδικασία για κάθε ανεξάρτητο αρχείο ήχου, ο επανυπολογισμός του σε κάθε εποχή κρίθηκε επιβαρυντικός για τη διαδικασία της εκπαίδευσης και άρα αποφεύχθηκε μέσω του προ-υπολογισμού και του cacheing των αποτελεσμάτων.

\subsection{Συνάρτηση Κόστους - VisemeAwareFrameLoss}
\label{LossFunction}

Για την εκπαίδευση του μοντέλου χρησιμοποιείται ένας προσαρμοσμένος στόχος μάθησης με τίτλο VisemeAwareFrameLoss. Η συνάρτηση αυτή συνδυάζει δύο όρους, έναν όρο cross-entropy στο επίπεδο των frames και έναν όρο Auxiliary CTC Loss. Για τεχνικές πληροφορίες περί των δύο αυτών όρων της αντικειμενικής συνάρτησης, ανατρέξτε στην ενότητα \ref{MachineLearningBackground}.

Για κάθε έγκυρο frame, υπολογίζεται ο όρος της cross-entropy ανάμεσα στην προβλεπόμενη κατανομή $q$ και την ομαλή κατανομή target $p$, με την εξομάλυνση που παρουσιάστηκε παραπάνω. Στον όρο αυτόν προστίθεται μία περαιτέρω ποινή με στόχο τα frames τα οποία έχουν κατηγοριοποιηθεί ως frames σιγής. Συγκεκριμένα, για frame με target $y$ και πιθανότητα πρόβλεψης του blank token $\varnothing$ ίση με $P(\varnothing)$, ο όρος ποινής $\phi$ υπολογίζεται ως 
\[
\phi = \lambda_i \cdot \mathbf{1}[y = \varnothing] \cdot (1 - P(\varnothing))
\]

όπου $\mathbf{1}[\cdot]$ είναι η δείκτρια συνάρτηση και $\lambda_i$ είναι παράμετρος που ορίζει τη μέγιστη τιμή του προστιθέμενου όρου με προεπιλεγμένη τιμή 0.5. Παρατηρώντας τον παραπάνω τύπο, συμπεραίνουμε ότι παίρνει τιμή 0 όταν το label του frame δεν αντιστοιχεί στη σιγή ή όταν αντιστοιχεί στη σιγή και το μοντέλο προβλέπει με απόλυτη βεβαιότητα το blank token. Σε περίπτωση που το blank token είναι το σωστό, όσο η πρόβλεψη του μοντέλου για το blank token γίνεται με λιγότερη αυτοπεποίθηση, τόσο αυξάνεται και ο όρος ποινής (γραμμικά), με μέγιστη τιμή το $\lambda_i$. Στόχος του όρου αυτού είναι να τιμωρούνται σε μεγαλύτερο βαθμό τα ``false positives'', δηλαδή η πρόβλεψη φωνημάτων σε frames στα οποία δεν υπάρχει κάποιο φώνημα. Ο συμμετρικός όρος, όρος δηλαδή που θα τιμωρούσε τα ``false negatives'', τις προβλέψεις της σιγής σε περίπτωση που στο frame αντιστοιχίζεται κάποιο φώνημα, δε συμπεριλαμβάνεται, καθώς κατά τις δοκιμές βρέθηκε ότι δεν επηρεάζει αισθητά το αποτέλεσμα.

Η ποινή προστίθεται στην τιμή που υπολογίστηκε από το cross-entropy και το άθροισμα αυτό πολλαπλασιάζεται με το confidence score της πρόβλεψης όπως παρουσιάστηκε προηγουμένως. Τέλος, υπολογίζεται ο μέσος όρος του παραπάνω για όλα τα έγκυρα frames, παράγοντας τον όρο $\mathcal{L}_f$ του batch. Πέραν του όρου αυτού, υπολογίζεται και η τυπική συνάρτηση κόστους CTC, η οποία ενισχύει το σήμα της σωστής αλληλουχίας φωνημάτων και είναι ανεξάρτητη από την οριοθέτηση των φωνημάτων σε επίπεδο frames. Ο όρος CTC $\mathcal{L}_c$ πολλαπλασιάζεται με τη σταθερά $\lambda_c$ με τυπική τιμή 0.3 και προστίθεται στον όρο $\mathcal{L}_f$ υπολογίζοντας την τελική τιμή της συνάρτησης κόστους:

\[
\mathcal{L} = \mathcal{L}_f + \lambda_c \cdot \mathcal{L}_c
\]

Στόχος του συνδυασμού των δύο συναρτήσεων κόστους είναι να δίνεται στο μοντέλο αυστηρή χρονική επίβλεψη (supervision), ειδικά όταν η πρόβλεψη του target γίνεται με αυτοπεποίθηση, ενώ ταυτόχρονα να υπάρχει το δεύτερο σήμα που δεν τιμωρεί υπερβολικά τις διαφωνίες στην οριοθέτηση αλλά εστιάζει στην επίτευξη της σωστής αλληλουχίας κατά την πρόβλεψη.

\subsection{Εκπαίδευση του Audiovisual Μοντέλου}
\label{AudiovisualTraining}

\paragraph{Modality Dropout}
\label{ModalityDropout}
Κατά την εκπαίδευση, εφαρμόζεται ολικό dropout του feature stream αυστηρά μίας εκ των δύο τροπικοτήτων με πιθανότητα $p_m$ (προεπιλεγμένη τιμή 0.15). Η επιλογή αυτή είναι συχνή σε περιπτώσεις όπου συνδυάζονται δύο διαφορετικές τροπικότητες και στοχεύει στην αύξηση της ευρωστίας του μοντέλου. Εξαναγκάζει τον εκάστοτε μηχανισμό συνδυασμού να μπορεί να χειρίζεται περιπτώσεις όπου μία ολόκληρη τροπικότητα είναι μη διαθέσιμη ή μη αξιόπιστη σε υψηλό βαθμό.

\paragraph{Κεφαλές Auxiliary CTC}
Παρέχεται η προαιρετική επιλογή (στην προεπιλογή ενεργοποιημένη) χρήσης auxiliary κεφαλών CTC, κατά την οποία προστίθενται δύο γραμμικές στρώσεις classifiers στις αναπαραστάσεις βίντεο και ήχου πριν τον συνδυασμό τους $F_v\sp{\prime}(t)$ και $F_a\sp{\prime}(t)$. Κάθε μία από αυτές παράγει τα $V+1$ logits ανά frame, τα οποία χρησιμοποιούνται κατά τον υπολογισμό της loss function (βλ. παρακάτω) και αγνοούνται στο βήμα του inference. Ο σκοπός της χρήσης των παραπάνω κεφαλών (της παραγωγής, δηλαδή, των auxiliary CTC όρων για βίντεο και ήχο) είναι προληπτικός, ώστε να περιοριστεί το πρόβλημα των vanishing gradients. Μία τροπικότητα η οποία συστηματικά λαμβάνει μικρότερο βάρος (επειδή παρουσιάζει, γενικά, στην κατανομή περισσότερο θόρυβο) μπορεί να λάμβανε σήμα vanishing gradients και άρα να μην προάγεται η εκμάθηση των αναπαραστάσεων της συγκεκριμένης τροπικότητας. Τότε, το μοντέλο θα έκλινε όλο και περισσότερο σε unimodal μοντέλο και δε θα υπήρχε η δυνατότητα να αξιοποιηθούν τα στοιχεία της γενικά πιο θορυβώδους τροπικότητας, όταν αυτά είναι αξιόπιστα.

\paragraph{Συνάρτηση Κόστους}
Η αντικειμενική συνάρτηση για το συγκεκριμένο μοντέλο είναι όμοια με την αντίστοιχη του visual μοντέλου που παρουσιάστηκε στην ενότητα \ref{LossFunction}, απαρτίζεται δηλαδή από τον όρο cross-entropy και τον όρο CTC, ενώ σε αυτούς προστίθενται και οι εκάστοτε όροι auxiliary CTC του βίντεο $\mathcal{L}_v$ και του ήχου $\mathcal{L}_a$ (ή ενός εκ των δύο, σε περίπτωση που επιλεχθεί η εκτέλεση με μοναδική τροπικότητα), πολλαπλασιασμένοι με τη σταθερά $\lambda_a$, με προεπιλεγμένη τιμή 0.1. Οι όροι αυτοί υπολογίζονται ακριβώς όπως ο όρος $\mathcal{L}_c$, χρησιμοποιώντας ωστόσο τα logits που παράγονται από τις κεφαλές auxiliary CTC και όχι τα logits της εξόδου μετά τον Sequence Encoder. Τελικά, η τιμή της συνάρτησης κόστους δίνεται ως:
\[
\mathcal{L} = \mathcal{L}_f + \lambda_c \cdot \mathcal{L}_c + \lambda_a \cdot (\mathcal{L}_v + \mathcal{L}_a)
\]

\subsection{Άλλες λεπτομέρειες για την εκπαίδευση}
\label{TrainingDetails}

Όσον αφορά την εκπαίδευση του μοντέλου, χρησιμοποιήθηκε AdamW optimizer (με weight decay χρησιμοποιώντας τυπική τιμή $10^{-4}$, με γραμμικό warmup του ρυθμού μάθησης (τυπικά για 5 εποχές) μέχρι τη μέγιστη τιμή ($3\cdot10^{-4}$). Χρησιμοποείται επίσης ReduceLROnPlateau με scheduling που βασίζεται στο validation token error rate (βλ. ενότητα \ref{Criteria}). Εφαρμόζεται gradient clipping στο 1.0, ενώ χρησιμοποιείται και μηχανισμός early stopping της εκπαίδευσης, αν το validation token error rate δε βελτιωθεί για τουλάχιστον $10^{-4}$ για έναν αριθμό εποχών (τυπικά 12). Για την αξιολόγηση στο test set, χρησιμοποιείται το μοντέλο με το χαμηλότερο validation TER, το οποίο δεν είναι αναγκαστικά αυτό της τελευταίας εποχής.

Να σημειωθεί ότι σε κάθε περίπτωση εφαρμόζεται gradient-norm clipping, το οποίο αποτελεί αναγκαία προσθήκη για την αποφυγή των explosive gradients ειδικά για την περίπτωση του BiLSTM, τα οποία οφείλονται στον πολλαπλασιασμό των recurrent gradients σε διάφορα βήματα του χρόνου. Από την άλλη, λόγω της αρχικής αστάθειας που παρουσιάζουν οι transformers ειδικά κατά τον υπολογισμό logits (τα αρχικά αποτελέσματα του attention είναι εσωτερικά γινόμενα τυχαίων προβολών query και key, με αποτέλεσμα οι προβλέψεις να είναι ιδιαίτερα θορυβώδεις πριν οι προβολές να έχουν μάθει κάτι ιδιαίτερο), χρησιμοποιείται linear warmup του ρυθμού μάθησης, κρατώντας τον χαμηλό όσο παράγονται μη αξιόπιστα gradients που ίσως οδηγούσαν το μοντέλο σε αστάθεια. Οι δύο αυτές προσαρμογές, αν και αντιτίθενται σε διαφορετικά προβλήματα των δύο μοντέλων, χρησιμοποιούνται και πάλι σε κάθε περίπτωση, ώστε η αρχιτεκτονική των μοντέλων να παραμένει σταθερή.

\subsection{Fine-tuning του LLM για την Πρόβλεψη Φράσεων}
\label{LLMFineTuning}

Για την πρόβλεψη φράσεων μέσω LLM (βλ. ενότητα \ref{PhrasePrediction}), έστω $S = (s_1, ..., s_M)$ η πραγματική πρόταση (αλληλουχία λέξεων $s_i, i=1,...,M$) που εκφέρεται από μία αλληλουχία φωνημάτων $P$. Εφαρμόζεται fine-tuning στο LLM χρησιμοποιώντας την παρακάτω Cross-Entropy Loss:

\[
\mathcal{L}_s = - \sum_{t=1}^{M} \ln p(s_t \mid P, s_{1:t-1})
\]
όπου $p(s_t|P, s_{1:t-1})$ είναι η πιθανότητα ορθής πρόβλεψης της επόμενης λέξης, με δεδομένο την αλληλουχία φωνημάτων και τις λέξεις που έχουν προβλεφθεί ως προηγούμενες.

Να σημειωθεί ότι η προηγούμενη προτεινόμενη μέθοδος δε δοκιμάζεται, διότι το μοντέλο δεν εφαρμόζεται σε περιβάλλοντα μη περιορισμένου λεξιλογίου.
